Different from the limited amount of data being produced from individual experiments, the large amount of materials data that are generated from PES sampling, high-throughput computations and stored in large-scale public materials database requires fundamentally new approaches to extract the underlying relationships amongst the data, in order to deepen our scientific understanding on materials' behaviours and facilitate new scientific discoveries. 

From the perspective of information theory, the process of data collection in computational materials science can be considered as a process of information/data compression. A typical calculation $\bm{R}\xrightarrow{\hat{\mathcal{H}}}\mathcal{A}$ that maps the three-dimensional input structure $\bm{R}$ into its corresponding properties $\mathcal{A}$ (\emph{e.g.} energy, charge density, electronic structures \emph{etc.}) via a Hamiltonian $\hat{\mathcal{H}}$ (\emph{e.g.} DFT), where the final data to be collected is typically a tuple $(\bm{R},\mathcal{A})$, will inevitably lose many detailed information about $\hat{\mathcal{H}}$, which encodes the intriguing knowledge on how many-particle interactions in a system that leads to the physical observable $\mathcal{A}$. To overcome this problem of information loss, building on the theory of probability discussed in \cref{sec:prob_theory_fundamentals}, ML provides a set of (light-weighted) statistical tools to help us infer the relationship of the data through discovering recurring patterns in the data. In the context of atomistic materials modelling, it learns an approximate mapping between the \emph{mathematical representation} $\mathcal{X}$ of the materials and their corresponding properties, through capturing the dominant (most probable) features in $\mathcal{X}$ that contributed towards $\mathcal{A}$. 

Formally, on the highest level, many ML models aim to learn the conditional distribution $p(y|\bm{R})$ that captures the likelihood of a structure $\bm{R}$ will possess some property or label $y$. Fundamentally, the underlying physical principle that enables the application of probabilistic (or statistical) models to describe materials is the belief that chemically and structurally similar materials should have similar properties, such that data sharing many commonalities should naturally aggregate together and occur with high probability. The most prominent example behind this belief is the establishment of the periodic table of elements that categorised elements into periods and groups, long before the birth of ML, or even quantum mechanics. In this regard, previous discussions on the theory of probability (\cref{sec:prob_theory_fundamentals}) will not only help to provides the foundation for understanding the following discussions on the mechanistic details of different ML/AI models, but as we hope, to promote the establishment of linkage between statistical mechanics, energy landscape sampling, and other aspects of materials modellings that have been discussed in this monograph. 

This section will focus on discussing the tasks achievable and models that are commonly applied in supervised and unsupervised machine learnings, as well as how these models can be trained and interpreted. For the time being, we shall assume that a proper mathematical \textit{descriptor} (also known as the \textit{feature vector}, which we shall use interchangeably in the following discussions) $\mathcal{X}$ that represent the underlying data (predominantly the atomistic structures) as the input to the ML models is already available, and shall return to the problem of how to construct $\mathcal{X}$ in \cref{sec:feature_eng}.
